% !PRINTED 25
\weberpage{25}

Die $a_i$ ändern sich, wenn man die $\xi$, $\eta$ einer Substitution (6)
unterwirft. Bildet man aber die Invarianten der Form (8) (nach
Algebra, Bd. I, §\,70), so erhält man Funktionen der Koeffizienten,
die auch diesen Substitutionen gegenüber invariant sind, die also
nicht zu den individuellen Formen $\varphi$, $\psi$, sondern zu dem
ganzen Büschel und also auch zu ihrem Durchschnitt, der Grundkurve,
gehören. Wir nennen sie \gesperrt{Invarianten der Grundkurve}. In bezug auf das
Formensystem $\varphi$, $\psi$ werden sie auch \gesperrt{Kombinanten} genannt.
Die Grundkurve hat also zwei Invarianten:
\weq{9}{\begin{aligned}
A &= a_2^2 - 3\,a_1a_3 + 12\,a_0a_4,\\
B &= 27\,a_1^2a_4 + 27\,a_0a_3^2 + 2\,a_2^3 - 72\,a_0a_2a_4 - 9\,a_1a_2a_3,
\end{aligned}}
aus denen man die \gesperrt{Diskriminante} $D$ nach der Formel
\weq{10}{27\,D = 4\,A^3 - B^2}
ableitet. In bezug auf die Koeffizienten von $\varphi$ und $\psi$ sind die
Invarianten $A$, $B$, $D$ von den Graden 8, 12, 24.

Das Formensystem $\varphi$, $\psi$ hat \gesperrt{zwei simultane quadratische
Kovarianten}, die man ebenso wie die Kovariante $C$ (in Algebra I,
§\,65) ableitet. Die erste von ihnen ist, wenn wir
$\psi_1 = \dfrac{1}{2}\dfrac{\partial \psi}{\partial x_1}, \dots$
setzen:
\weq{11}{\Phi = \begin{vmatrix}
a_{11}, & a_{12}, & a_{13}, & a_{14}, & \psi_1\\
a_{21}, & a_{22}, & a_{23}, & a_{24}, & \psi_2\\
a_{31}, & a_{32}, & a_{33}, & a_{34}, & \psi_3\\
a_{41}, & a_{42}, & a_{43}, & a_{44}, & \psi_4\\
\psi_1, & \psi_2, & \psi_3, & \psi_4, & 0
\end{vmatrix}}
und die zweite $\Psi$ erhält man daraus, indem man $a_{ik}$ mit $b_{ik}$
und $\psi$ mit $\varphi$ vertauscht.

Die Gleichung $\Phi = 0$ drückt eine Fläche zweiten Grades aus (die
nicht zum Büschel gehört), die der geometrische Ort der Punkte $x$
ist, deren Polaren in bezug auf $\psi$ die Fläche $\varphi$ berühren.

Wir nehmen jetzt an, daß die Diskriminante $D$ von Null verschieden
ist. Dann lassen sich die beiden Funktionen $\varphi$, $\psi$
simultan in die Summe von vier Quadraten transformieren:
\weq{12}{\begin{aligned}
\varphi &= \alpha_1y_1^2 + \alpha_2y_2^2 + \alpha_3y_3^2 + \alpha_4y_4^2,\\
\psi &= \beta_1y_1^2 + \beta_2y_2^2 + \beta_3y_3^2 + \beta_4y_4^2,
\end{aligned}}
worin die $y_i$ lineare Funktionen der $x_i$ sind. Dies ist die
\gesperrt{kanonische Form} des Funktionenpaares. Es wird dann:
\weq{13}{f(\xi,\eta) = (\xi\alpha_1+\eta\beta_1)(\xi\alpha_2+\eta\beta_2)(\xi\alpha_3+\eta\beta_3)(\xi\alpha_4+\eta\beta_4),}
%CONT>>
